\subsection{Sprint 1}

Nesta sprint, fiquei responsável pelo desenvolvimento do componente de carrossel (US-03) e pela implementação da tela de introdução (US-03). Porém, antes de iniciar o desenvolvimento das \textit{tasks}, tomei a iniciativa de criar um tema customizado para a aplicação, com o objetivo de garantir que todas as telas seguissem um padrão visual consistente e harmonioso. A definição de cores e fontes no tema, bem como sua acessibilidade em todo o código, contribuíram para uma identidade visual clara e coesa em cada uma das telas. Isso facilitou não apenas o desenvolvimento estético e visual com base no protótipo das telas criado na sprint 0, mas também a manutenção do projeto.

Em seguida, foquei na criação do componente de carrossel, para o qual optei por utilizar a biblioteca Swiper. A escolha dessa ferramenta permitiu o desenvolvimento de um componente altamente customizável e que pode ser reutilizado em outras partes da aplicação. Isso facilitará significativamente o desenvolvimento das próximas \ac{us} que também incluirão componentes de carrossel semelhantes ou que exijam navegação por conteúdo visual. Durante o desenvolvimento da tela de introdução, colaborei diretamente com minha colega de AGES I, assegurando que o carrossel fosse integrado corretamente à nova interface. Após essa etapa, passamos por uma revisão de código, onde ajustes foram sugeridos pela equipe. Atendi às solicitações e fiz as modificações necessárias, adicionando responsividade à tela e implementando os testes adequados. Ao longo da sprint, ofereci suporte à minha dupla, explicando todo o processo de desenvolvimento e esclarecendo dúvidas quando surgiam. Além disso, auxiliei meus colegas sempre que necessário e contribuí para outras tarefas fora do escopo da sprint, mas igualmente importantes, como a implementação de testes unitários e de componentes.

 Assim, também tomei a iniciativa de implementar um \textit{customRender} para aprimorar o ambiente de testes, o que permitiu a criação um cenário mais robusto, capaz de replicar de forma mais fiel a interação de um usuário com os componentes. Isso garantiu que os testes realizados fossem mais confiáveis e eficientes. Ao final, utilizei o \textit{customRender} no ambiente de testes da tela de introdução e no componente de carrossel, resultando em uma boa cobertura de testes para ambos.

O principal desafio técnico desta sprint foi relacionado à configuração de um framework de testes adequado para a aplicação. Inicialmente, planejei utilizar o Jest em conjunto com a React Testing Library. No entanto, surgiram diversos erros complexos durante a implementação, difíceis de diagnosticar e resolver. Percebendo que estava gastando muito tempo sem encontrar uma solução satisfatória, busquei ajuda dos meus colegas, e juntos decidimos migrar para o Vitest, uma alternativa mais leve e eficaz para o nosso contexto. A transição para o Vitest foi uma escolha estratégica, permitindo solucionar mais facilmente os problemas iniciais de configuração e resultando em uma implementação tranquila e produtiva de testes unitários e de componentes.

A experiência desta sprint me ensinou a importância de saber quando pedir ajuda e a necessidade de ser flexível nas abordagens adotadas. Às vezes, abandonar uma ideia inicial que não estava funcionando pode abrir caminho para soluções mais eficazes e melhor adaptadas às necessidades do projeto. A colaboração e a troca de conhecimentos com os membros da equipe foram fundamentais para que pudéssemos solucionar rapidamente os problemas técnicos que surgiram, evitando gargalos e outros impedimentos.

Com o encerramento desta sprint, o próximo passo será direcionar meus os esforços para o desenvolvimento das próximas user stories que me forem atribuídas. Continuaremos avançando de acordo com as prioridades definidas pela equipe, sempre buscando melhorar o fluxo de trabalho e a qualidade das entregas nas próximas etapas do projeto.